\documentclass{resume}

\newcommand{\en}[1]{}
\newcommand{\zh}[1]{#1}
\geometry{left=0.45in,right=0.45in,top=0.14in,bottom=0.20in}
\setlist[itemize]{itemsep=0.15ex,topsep=0.15em,leftmargin=1.25pc}

\zh{
  \usepackage{xeCJK}
  % SimSun has no real bold/italic. FakeBold on upright makes body slightly
  % heavier; larger FakeBold on bold keeps \textbf distinct.
  \newcommand{\resumeCJKFaces}[1]{%
    \setCJKmainfont{#1}[
      UprightFont=*, BoldFont=*, ItalicFont=*, BoldItalicFont=*,
      UprightFeatures={FakeBold=1.8},
      BoldFeatures={FakeBold=4.6},
      ItalicFeatures={FakeBold=1.8, FakeSlant=0.167},
      BoldItalicFeatures={FakeBold=4.6, FakeSlant=0.167}]%
    \setCJKsansfont{#1}[
      UprightFont=*, BoldFont=*, ItalicFont=*, BoldItalicFont=*,
      UprightFeatures={FakeBold=1.8},
      BoldFeatures={FakeBold=4.6},
      ItalicFeatures={FakeBold=1.8, FakeSlant=0.167},
      BoldItalicFeatures={FakeBold=4.6, FakeSlant=0.167}]%
    \setCJKmonofont{#1}[
      UprightFont=*, BoldFont=*, ItalicFont=*, BoldItalicFont=*,
      UprightFeatures={FakeBold=1.8},
      BoldFeatures={FakeBold=4.6},
      ItalicFeatures={FakeBold=1.8, FakeSlant=0.167},
      BoldItalicFeatures={FakeBold=4.6, FakeSlant=0.167}]%
  }
  \IfFontExistsTF{SimSun}{
    \resumeCJKFaces{SimSun}
  }{
    \IfFontExistsTF{Songti SC}{
      \resumeCJKFaces{Songti SC}
    }{
      \IfFontExistsTF{Noto Serif CJK SC}{
        \setCJKmainfont{Noto Serif CJK SC}[UprightFont=* Medium, BoldFont=* Bold]
        \setCJKsansfont{Noto Serif CJK SC}[UprightFont=* Medium, BoldFont=* Bold]
        \setCJKmonofont{Noto Serif CJK SC}[UprightFont=* Medium, BoldFont=* Bold]
      }{
        \setCJKmainfont{FandolSong-Regular}[BoldFont=FandolSong-Bold]
        \setCJKsansfont{FandolSong-Regular}[BoldFont=FandolSong-Bold]
        \setCJKmonofont{FandolSong-Regular}[BoldFont=FandolSong-Bold]
      }
    }
  }
  \setmainfont{Latin Modern Roman}[
    UprightFeatures={FakeBold=1.4},
    ItalicFeatures={FakeBold=1.4}]
  \setmonofont{Latin Modern Mono}[
    UprightFeatures={FakeBold=1.4},
    ItalicFeatures={FakeBold=1.4}]
}

\begin{document}
\pagestyle{empty}
\fontsize{8.6pt}{13.0pt}\selectfont
\emergencystretch=1.5em
\vspace*{-0.8ex}

\noindent
\begin{minipage}[t]{0.74\textwidth}
  {\LARGE\scshape \en{Qingliu Wu}\zh{吴清柳}\par}
  \vspace{0.6ex}
  {\ttfamily\normalsize
    \phone{17353675760} \textperiodcentered\
    \email{wql@bupt.edu.cn} \textperiodcentered\
    \github[micuks]{https://github.com/micuks}\par
  }
\end{minipage}\hfill
\makebox[0pt][r]{\raisebox{-7ex}[0pt][0pt]{\includegraphics[width=2.2cm,clip,trim=0 57 0 0]{selfie.png}}}
\vspace{-0.2ex}

\section{\en{Education}\zh{教育经历}}
\en{\datedsubsection{\textbf{Beijing University of Posts and Telecommunications}, Master Student}{09/2023 -- Present}}
\zh{\datedsubsection{\textbf{北京邮电大学}, 硕士研究生}{2024/09 -- 2027/06}}
\begin{itemize}[parsep=0.25ex]
	\item \zh{计算机科学与技术（学硕），获\textbf{2024--2025 年度国家奖学金}；IDMG 实验室，师从邵蓥侠老师，方向 \textbf{AI4DB}}
\end{itemize}
\en{\datedsubsection{\textbf{Beijing University of Posts and Telecommunications}, Bachelor of Engineering}{09/2020 -- 06/2024}}
\zh{\datedsubsection{\textbf{北京邮电大学}, 工学学士}{2020/09 -- 2024/06}}
\begin{itemize}[parsep=0.25ex]
	\item \zh{计算机科学与技术（计算机学院实验班），综合成绩 91.91，排名 29/422（前 6.56\%）}
\end{itemize}

\section{\zh{实习经历}}

\zh{\role{Qwen 系列大模型训练优化 · FP8 低精度训练 / Megatron}{2026.10 -- 至今}}
\zh{{\small\itshape 快手 · 算法引擎部 训练引擎组 | AI Infra}\par}
\begin{itemize}[parsep=0.25ex]
	\item \zh{\textbf{FP8 低精度训练显存优化：}面向 Qwen3.6-35B-A3B VLM 的 32 卡 H100、131K 长序列训练，定位 Megatron FP8 训练相比 BF16 多出的 2.5 GiB 几乎全部来自优化器状态；据此定位并实现 CE 权重梯度直接累加进主梯度、attention 低精度 row/col buffer在最后一个消费者后提前释放且权重列副本有界缓存复用、q\_layernorm 与 RoPE 改为重算、GDN 复算过程减少一次前向等优化，均验证逐位一致，峰值显存下降 8.1\%（51,872 $\to$ 47,659 MiB），同时训练单步耗时基本持平。}
	\item \zh{\textbf{MoE 低精度算子优化：}将 CuTe DSL MoE 融合算子（MegaMoE）接入 Megatron 的 MoE 数据流，用 Triton 把 FP8 专家权重准备中逐权重反量化、拼接、gate/up 行交错、再量化的 203 个碎 kernel 合并为批量算子，kernel launch 数降低约 98\%、端到端步时 $-$4.32\%；配合调度优化与缓存复用，该路线单步由 11,280 降至 10,420 ms（$-$7.6\%），反超 BF16 基线 1.6\%。}
	\item \zh{\textbf{流水线与并行训练支撑：}定位流水线末级因额外输出头负载叠加 P2P 后的全设备同步而放大的串行等待，去掉该同步后 32 卡步时降低 7.1\%、跨 stage 通信等待由 3.205 s 降至 1.409 s；支撑 Qwen3.8-27B 的 Megatron-FSDP H200 训练，并实现重叠参数聚合与视觉塔同步。}
\end{itemize}

\zh{\role{多模态大模型 Keye 预训练 · Megatron/DeepSpeed 框架优化}{2026.05 -- 2026.10}}
\zh{{\small\itshape 快手 · 基础大模型与应用部 | AI Infra}\par}
\begin{itemize}[parsep=0.25ex]
	\item \zh{\textbf{实现长序列训练架构：}为 Causal-ViT 实现 Context Parallel，为 GDN 实现 Ulysses All-to-All、为 MoDA 时序注意力实现跨分片 halo 通信；支持图片--视频混训，实现两者的混训的分片与损失归约，解除 CP 路径 batch size$=$1 限制，将 CP8 单样本训练长度从 131K 扩展至 2M visual tokens。}
	\item \zh{\textbf{支撑千卡训练稳定性：}支撑 Stage-2 千卡长视频预训练（CP8$\times$DP128、GBS 2,048、单样本最高 2,048 帧/约 2M visual tokens）；定位 CP8 短视频空分片漏发反向 All-Reduce 导致 collective 序列失配与训练 hang，通过计算图锚点统一各 rank 通信序列，修复后在千卡训练中无复现。}
	\item \zh{\textbf{MoDA 计算性能优化：}将 launch-bound、GPU 空转约 80\% 的逐帧 Python 循环重构为跨帧批处理，支持 batch size$>$1，单层/端到端提速 5.4$\times$/1.28$\times$；进一步设计 chunked packed MoDA，按 QKV、gate、attention 与 output 分阶段切分中间张量配合反向传播重算，在显存峰值小幅增加，仍旧支持 2M 序列训练的同时，单层提速 3.36x，并完成 BF16 单层 MoDA 前反梯度对齐，真实训练 eval 表现无下降。}
	\item \zh{\textbf{分布式优化器：}定位 Muon 与 ZeRO-1/2 reduce-scatter 的不兼容根因——分区参数矩阵缺少 Newton-Schulz 所需的完整归约梯度；向 DeepSpeed 上游提交切片定向传播与 owner rank 冗余更新修复，\textbf{PR\#8141 已合入 master}。}
	\item \zh{\textbf{工程正确性：}保障 DeepSeek-V4-Flash 预训练数值正确性：定位 fused RoPE 触发 BF16 训练 NaN 的根因，提交含最小复现的上游 issue（Megatron-LM\#5318）并推进修复。}
\end{itemize}

\zh{\role{Flink 大状态计算调优}{2026.03 -- 2026.05}}
\zh{{\small\itshape 华为 · 计算产品线 | AI 工程师}\par}
\begin{itemize}[parsep=0.25ex]
	\item \zh{\textbf{状态访问链路优化：}面向 Flink 1.16.3 + RocksDB 大状态场景，构建 Value/Map 双层堆内缓存与 MapState EMPTY/SINGLE 快照；进一步以 Mailbox 攒批驱动 MultiGet 异步预取和同 key 本地预聚合，并在 MemTable/SST 前加入 whole-key Bloom miss gate，系统性减少序列化、JNI 往返与无效读取。}
	\item \zh{\textbf{在鲲鹏平台取得显著性能收益：}在亿级数据量下对 Nexmark 评测中，较 RocksDB StateBackend 的逐 Query 平均吞吐在鲲鹏平台提升 50\%，并通过 CDC 数据一致性检验。}
\end{itemize}

\section{\zh{项目经历}}
\zh{\role{SQLForge：执行反馈驱动的 Text-to-SQL Agent 后训练\textnormal{\normalfont\itshape{ | QLoRA / GRPO / vLLM}}}{2026.01 -- 2026.05}}
\begin{itemize}[parsep=0.25ex]
	\item \zh{\textbf{执行反馈后训练与交互式自纠错：}构建 Agent 后训练管线——生成多 teacher ReAct 轨迹并用真实执行结果回填 Observation，QLoRA SFT 后接自研 GRPO 循环与格式/合法性/结构/部分执行/最终正确性稠密奖励；推理侧设计 Direct SQL + 执行失败后 ReAct 多轮交互，在沙箱中执行草稿 SQL 并依据 Observation 修正，Spider 执行正确率 78.6\%，较 ReAct SFT direct-only 的 71.7\% 提升 6.9pp（Hard 难度 +13.4\%）。}
\end{itemize}


\section{\zh{科研经历}}
\zh{\role{DBPecker: 分布式数据库复合异常诊断系统}{PVLDB 2025 · 共同一作}}
\begin{itemize}[parsep=0.25ex]
	\item \zh{\textbf{复合异常诊断：}构建跨节点 compound anomaly graph，以相关性加权 PageRank 融合节点影响力与多标签异常概率，输出按影响排序的 $\langle$node, anomaly$\rangle$ 序列；通过结构因果模型干预识别跨 OS/数据库层根因指标，并在 OceanBase 集群以故障注入完成端到端验证。}
\end{itemize}

\section{\en{Skills}\zh{技能}}
\begin{itemize}[parsep=0.25ex]
	\item \zh{\textbf{大模型系统：}Megatron-LM、DeepSpeed、Transformer Engine FP8、Triton、PyTorch、vLLM、SFT/GRPO/ReAct；千卡预训练性能、显存与 kernel 优化；}
	\item \zh{\textbf{编程与工具：}熟练使用 C++/Python 与 Claude Code、Codex；了解 Java/Go/Rust；熟悉 Linux、Git、Docker、CMake；}
	\item \zh{\textbf{英语与荣誉：}雅思 7.5（阅读/听力 8.5）、六级 594；中国机器人及人工智能大赛百度 Apollo 自动驾驶仿真赛全国一等奖。}
\end{itemize}

\end{document}
